% =============================================================================
% §4 Discrimination — identification (M3 / R3). β>0 everywhere; COL cohort-
% dominant + small exogenous period; CRI cohort-leaning; 2c-i individual
% confirmation; explicit identification-strategy statement; ecological caveats.
% =============================================================================
\section{Discriminating the Mechanism: APC Curvature and the Sign of State
Dependence}
\label{sec:discrimination}

\subsection{Identification strategy, stated plainly}
\label{sec:idstrategy}

Two quantities identify the discrimination; both were fixed, with their decision
rule, before the estimates were run.

\emph{First, where the change localizes.} We build birth-year pseudo-cohorts
from the repeated cross-sections and decompose the married (and cohabiting)
share into age, period, and cohort components. Because the linear APC trends
are not separately identified, we rest nothing on them: only
\emph{curvature} (identified second differences) enters the verdict --- where
does the \emph{acceleration} sit, on cohort lines or on period lines? Age-band
fixed effects and within-band estimation keep Lexis artifacts out of the
curvature.

\emph{Second, the sign of state dependence.} For the cascade hypothesis the
decisive quantity is the coefficient $\beta$ on the \emph{lagged reference-group
not-married share} in a within-pseudo-cohort regression of the change in
marriage outcomes. An amplifying cascade requires $\beta < 0$ and significant:
decline must be faster where non-marriage is already prevalent. A positive
$\beta$ --- decline \emph{smaller} where the reference share is high --- is the
opposite of the cascade signature: stabilizing, mean-reverting dynamics. To
guard against mechanical correlation, the reference group is varied across
specifications: the adjacent \emph{younger} cohort (primary), the adjacent
\emph{older} cohort (clean of own-cohort reflection), the whole 20--39
population (robustness), and the cohort's own lag (comparator). The
pre-registered rule required an amplifying $\beta$ under the primary and at
least one clean specification, in \emph{both} countries, for
H$_{\text{cascade}}$ to stand.

\subsection{The verdict: stabilizing, not self-reinforcing}
\label{sec:verdict}

Table~\ref{tab:beta} reports the full $\beta$ matrix at lag one, where the
signal concentrates (lag-two estimates collapse toward zero).

\begin{table}[t]
\centering
\caption{State-dependence coefficient $\beta$ on the lagged reference
not-married share (lag 1). Amplifying cascade requires $\beta<0$.}
\label{tab:beta}
\begin{tabular}{lcc}
\toprule
Reference-group specification & $\beta$ Costa Rica & $\beta$ Colombia \\
\midrule
Adjacent younger cohort (primary) & $+0.17$ & $+1.33$ \\
Adjacent older cohort (clean)     & $+0.16$ & $+0.61$ \\
Population 20--39 (robustness)    & $+0.29$ & $+0.97$ \\
Own lag (comparator)              & $+0.36$ & $+0.98$ \\
\bottomrule
\end{tabular}
\end{table}

$\beta > 0$ in every specification, in both countries: \emph{the cascade
signature is absent}. Within-cohort marriage decline is \emph{smaller}, not
larger, where the reference not-married share is already high, and the
pre-registered H$_{\text{cascade}}$ condition fails comprehensively.
Descriptively the estimated dynamics are mean-reverting; we use
``stabilizing'' in that descriptive sense only, without attaching a
behavioral mechanism to the positive sign --- the design licenses the
rejection of the amplifying signature, not an interpretation of its opposite.

The localization completes the picture. In Colombia, curvature localizes on
\emph{cohort} lines (married-share curvature RMS 0.007 on period lines versus
0.014 on cohort lines), with a statistically unambiguous period component
(F-test $p \approx 3\times10^{-6}$) that is \emph{small} relative to the cohort
step-down --- and, crucially, non-reflexive: because $\beta$ is not amplifying,
that period component carries an \emph{exogenous-shock} signature, not
endogenous feedback. In Costa Rica the period effect is not statistically
present at all ($p = 0.18$), leaving a cohort-replacement reading. The
structure, in both countries, is
H$_{\text{cohort}}$ --- dominant cohort replacement --- with a secondary
H$_{\text{shock}}$ overlay in Colombia. Figure~\ref{fig:apc} shows the
pseudo-cohort lines and the identified curvature; Figure~\ref{fig:beta} the
$\beta$ estimates with confidence intervals by specification.

\begin{figure}[t]
\centering
\begin{subfigure}[b]{0.48\textwidth}
  \includegraphics[width=\textwidth]{fig_cohortlines_COL.pdf}
  \caption{Colombia: pseudo-cohort lines}
\end{subfigure}
\hfill
\begin{subfigure}[b]{0.48\textwidth}
  \includegraphics[width=\textwidth]{fig_apc_curvature_COL.pdf}
  \caption{Colombia: identified APC curvature}
\end{subfigure}
\caption{Localization of the change (Colombia; Costa Rica in the online
appendix). Acceleration sits on cohort lines; the period overlay is real but
small.}
\label{fig:apc}
\end{figure}

\begin{figure}[t]
\centering
\begin{subfigure}[b]{0.48\textwidth}
  \includegraphics[width=\textwidth]{fig_beta_ci_CRI.pdf}
  \caption{Costa Rica}
\end{subfigure}
\hfill
\begin{subfigure}[b]{0.48\textwidth}
  \includegraphics[width=\textwidth]{fig_beta_ci_COL.pdf}
  \caption{Colombia}
\end{subfigure}
\caption{State-dependence $\beta$ by reference-group specification, with
confidence intervals. Positive everywhere: stabilizing, not cascading.}
\label{fig:beta}
\end{figure}

\subsection{The individual-level confirmation}
\label{sec:individual}

The estimates above are cell-level: on repeated cross-sections, the
state-dependence test is a within-pseudo-cohort \emph{cell} regression, and an
ecological objection is available --- individual dynamics could differ from
cell dynamics. For Colombia, where person-level microdata permit it, we re-ran
the test at the individual-outcome level (a multinomial model of current union
status on the lagged reference share, person weights, within cohort--age
cells). The individual-level estimates confirm the cell-level verdict --- no
amplifying state dependence --- discharging the ecological objection for the
country where the period overlay made it most pointed. What repeated
cross-sections cannot deliver at any level is \emph{true individual
transitions}; that requires panel or event-history data (the Colombian ENDS
program), which we return to in Section~\ref{sec:limitations}.

Two caveats bound the verdict rather than undermine it. The 20--39 estimation
floor excludes the 15--19 entry margin --- a cascade igniting among teenagers
would be under-observed here; and both countries are high-cohabitation regimes,
so the non-cascade verdict does not automatically travel to marriage-dominant
systems (that is the out-of-sample prediction of
Section~\ref{sec:interpretation}, not an assumption).
